\section{预实验与后续验证方案}
\label{sec:experiments}
本节将已执行结果与待执行实验分开。已有结果来自冻结提交\texttt{061ecf5ddea0a59dcf932b5038ad646b7af454a4}及其原始收据\ccite{V7Native}；本轮仅核对、重排和复算已记录的统计量，没有生成新气体世界、训练新滤波器或启动机器人闭环。软件可运行、读出与旧实现一致以及定位性能改善分别评价。

\subsection{数据来源、划分和读出资格}
原生归档为House01中\texttt{1,3-2,4\_fast}设置的两个释放位置，每个位置4个独立实现，共8个世界。原始逐世界1803帧，总计14424帧，覆盖0--999.502991\,s，并逐帧核对SHA256。其环境属于VGR/GADEN研究数据背景\ccite{Ojeda2023VGR}，归档名称中的\texttt{s2x}表示移源世界，不是枫树岭S2湖岸数据。源A和源B的三维坐标分别为$(-0.60,1.95,0.40)$\,m与$(-0.40,-2.90,-0.30)$\,m，均经原三维地图核查处于自由体素。两源同时改变水平位置和高度，不能将响应差解释为纯水平位置干预。

各源的r01/r02用于TRAIN，r03/r04用于DEV，因此训练和开发均为4个世界。两侧探针沿用此前冻结的12/17号，每个探针实际为4个原生格中心浓度的均值；两个采样池中心间距约7.003\,m，不能称为理想无体积单点PID。窗口在读取浓度前固定为50--70\,s，实际含33帧，时间为50.2998047--69.4995117\,s，相邻帧约0.6\,s；没有插值补造高频运动。每个世界贡献一个窗口，33帧不是33个独立实验。

原V7核库要求最大间隔0.25\,s，不能直接用于上述归档。原执行记录在开发评分前另行冻结粗采样适配：最大间隔0.61\,s，延迟1.8/3/6\,s，衰减1.5\,s，保留镜像算子及选择规则。该适配只允许解释粗采样、理想零延迟浓度读出的结果，真实传感器响应标定标记仍为未通过。它不代表恢复了动物尺度的快速气味运动，也不代表双无人机动力学或旋翼干扰已经验证。

原执行回归包括30个探针、10个时刻和4个训练世界的1200项浓度查询，全部通过，其中485项非零，最大绝对差$3.5763\times10^{-7}$\,ppm。软件测试最终71/71通过；20项C++/Python模型推理对照的符号及候选后验一致，最大特征差$2.22\times10^{-16}$。这些结果支持归档读出与实现一致性，不证明真实物理分布准确、概率校准或气源定位有效。

\subsection{固定观测下的完整普通对照}
所有方法使用相同TRAIN/DEV世界、相同窗口和候选支持，模型及预处理在DEV评分前冻结。表\ref{tab:native}保留原始\texttt{R2\_RESULT.json}汇总中的全部11个比较项。Gaussian头按有限训练窗口拟合；完整raw-joint同时提供向量量化（VQ）和Gaussian对照，以避免只用弱离散基线衬托方法。这里的“完整”指该窗口中实际保存的两侧样本，不包含窗口外的气体全场。

\begin{table*}[t]
\centering\small
\caption{两已标定源、4个独立开发世界上的冻结原生浓度读出结果。全部为固定观测离线对照；没有原生PMFS或自主双机成绩。}
\label{tab:native}
\begin{tabularx}{\textwidth}{Y c c r p{4.3cm}}
\toprule
方法 & 严格正确/4 & 并列数 & 原记录MAP均误差/m & 结果解释\\
\midrule
普通单点时间均值 + Gaussian & 4 & 0 & 0 & 两候选的当前开发分类已饱和\\
普通单点完整时间窗 + Gaussian & 4 & 0 & 0 & 单侧时间序列的普通控制\\
普通双点时间均值 + Gaussian & 4 & 0 & 0 & 未比单点控制增加判对样本\\
完整双点窗口 raw-joint VQ & 4 & 0 & 0 & VQ码本因4个训练块缩减为4\\
完整双点窗口 + Gaussian & 4 & 0 & 0 & 强普通联合对照\\
母结构best-one + 原直方图 & 2 & 2 & $0^{\dagger}$ & 实际选中静态梯度，非母文完整学习模型\\
母结构所选特征 + Gaussian & 4 & 0 & 0 & 同一受限静态特征，更换统计头后全部判对\\
固定双通道 + 原直方图 & 2 & 2 & $0^{\dagger}$ & 两个样本依赖固定顺序破并列\\
共同强度even + 原直方图 & 2 & 0 & 2.427061 & 与空通道选择后的M1一致\\
M1源对覆盖 + 原直方图 & 2 & 0 & 2.427061 & 动态镜像通道选择数为0\\
M1所选特征 + Gaussian头 & 1 & 0 & 3.640591 & 更换头部未恢复普通对照表现\\
\bottomrule
\multicolumn{5}{p{\dimexpr\textwidth-2\tabcolsep\relax}}{\footnotesize $\dagger$ 原记录按固定候选顺序破并列，恰好得到0\,m；严格判对仍为2/4。该0不作为无歧义定位成功。其余0也只是当前两个已标定候选的离线分类零误差，不表示连续未知源达到零米精度。}\\
\end{tabularx}
\end{table*}

M1在两个源A开发世界中均选择源B，在两个源B世界中判对，平均平面误差为2.427061\,m。共同强度的训练JS约0.242586，超过覆盖阈值0.04；唯一源对已被截断目标判为覆盖，贪心第一轮增益为0，最终选择集合为空。该行为由训练收据和导出模型直接确认，M1成绩也与共同强度even对照一致。母结构受限特征在Gaussian头下为4/4，进一步表明表征与统计头的作用需要分别控制，但它仍不是原母文的可学习网络。由于动态镜像通道没有进入M1评分链，本结果检验的是当前选择与量化方案，不能解释为动态镜像核已被充分使用后仍无效。

普通单点时间均值已经全部判对，说明当前任务没有建立双点或动态关系的必要性；M1的失败也不能简单归结为双机几何不佳。相同测量中存在普通方法能够利用的源差，而当前压缩未稳定保留。该4世界小样本不支持稳定泛化、显著性或真实概率校准主张，更不能在基线误差为零的任务上计算相对误差改善10\%。当前有限条件表仅覆盖两个已标定位置，未验证未见源泛化。

此前静态pilot和解析控制后端仅提供不同层次的辅助记录。静态pilot的普通raw-joint为40/40、静态M1为31/40，但40张快照来自4个开发世界；解析控制矩阵包含156个episode，其普通强基线与完整方法最终端点均饱和\ccite{V7Archive}。它们不是本节原生动态对照的追加独立样本，也不作为原生双机闭环收益。当前M1未通过独立增益门，M2/M3正式原生主动实验按原合同未启动，保留原始首次结果。

\subsection{需要逐项回答的实验问题}
后续实验围绕三个独立命题设置，而不以整套方法胜过一个弱基线替代模块验证。命题H1为：在训练预算和模型容量相同的条件下，源判别镜像通道能否较普通完整联合观测及母结构适配获得稳定源排序、校准或样本效率收益。命题H2为：在两机资源相同的条件下，边界记忆结合响应差几何能否比普通重捕获与无记忆判别几何取得更有辨识力的测量。命题H3为：在相同原始流、运动约束和后验更新下，完整源覆盖与新后缀审计能否降低错误高置信持续时间、改善恢复，而不增加正常任务错误。

\begin{table*}[t]
\centering\small
\caption{正式实验矩阵与执行状态。待执行项使用明确占位，不预填优势曲线。}
\label{tab:plan}
\begin{tabularx}{\textwidth}{p{1.45cm}Y Y p{1.9cm}}
\toprule
实验层次 & 比较对象 & 主要控制与回答的问题 & 当前状态\\
\midrule
固定观测H1 & 普通均值、raw-joint、母核库best-one、固定等数量通道、M1、无共模/无镜像/无源对目标 & 相同同步窗口、同训练块数、相同头部与总平滑质量；时间打乱、左右交换、传感器延迟扰动分开 & 两源开发完成；跨源结果待补\\
母文迁移 & PNAS原可学习结构 + 源条件头、固定核库best-one & 区分原母网络学习能力、必要任务适配和本文源对选择；训练预算相同 & 精确迁移待实现/待验证\\
几何H2 & 普通边界返回、无记忆响应差、仅返回区域、完整M2 & 相同双机、候选动作及代价；分别检查间距、方向、返回记忆和区域外复测 & 原生结果待补\\
决策H3 & 联合EIG、直接成本Chernoff、无覆盖、无审计反馈、完整M3 & 相同M1/M2与全原始流；正常先验和预设错误先验分别测试 & 原生结果待补\\
完整定位 & 原生PMFS-1U、透明适配PMFS-2U、raw-joint+EIG、各模块消融、完整方法 & 固定轨迹与自主闭环分开；设备数、样本数、墙钟与两机总时间分别核算 & 全域原生成绩待补\\
迁移与实用性 & 未见源、风、独立传播实现、House01/02/03、合格S2 & 未见源需先具备参数化条件预测接口；真实可用局地风与全真风oracle分开 & 待执行；S2资格未完成\\
\bottomrule
\end{tabularx}
\end{table*}

PMFS-1U保留原生算法和固定设置，作为设备资源参照；PMFS-2U是需公开说明的双输入适配，不冒称已有官方双机实现。主算法比较应在相同两机输入上进行，另报单机资源消融。M1的母方法对照至少包含同核库best-one与原可学习结构的任务适配，两者不能混称精确母方法复现。普通完整联合模型可包含VQ、正则Gaussian及适当容量的时间模型；模型选择只使用独立训练/验证资料，不能反复查看最终测试后挑选有利基线。

H1首先使用同观测固定轨迹，比较表示对同一批证据的利用。H2/H3和完整方法随后在原生GADEN-RT、相同合法起点和预算下运行自主闭环，使策略真正改变后续位置。当前代码在行进期读取气体事件供M2使用、在驻留期更新源概率，因此公平对照必须给所有方法相同可用原始流，或对这一额外信息单独消融；仅把飞行时间相等不能消除信息权限差异。真实运动、稳定、恢复和采样均计入成本，不把两个同时工作的传感器计为一个。

\subsection{划分、指标与统计判据}
独立性以完整传播世界为基本单位。训练、调参验证和最终测试按源位置、风条件和实现标识预先分开，不能把同一世界的相邻窗口分到两侧充当独立样本。House01/02/03分别报告，不把源数、风数和实现数相乘成没有实际运行的全交叉样本量。未见源验证还要求源坐标到条件响应的参数化模型或经资格审查的在线前向接口；当前按源编号标定的有限表不具备该能力，补齐接口之前该项保持待执行。

源级结果至少包括严格Top-1、MAP距离、Top-5\%位置读出、真源平均秩及源对分离方向。定义$\mathcal M_t=\operatorname*{arg\,max}_{g}P_t(g)$；严格Top-1只有在真源单元为唯一最大项时记为正确，并列单独报告。Top-5\%沿用固定官方O24实现的空间读出与误差计算，正式测试前须一并冻结其代码版本、合法单元域、取整、权重及并列规则。本文保留B4原CSV值及已有独立复算值，不在尚未逐项核验该实现细节时自行添加同名公式。若另设概率加权位置指标，应另名、另报，不能替换或混称官方Top-5\%。

后验可靠性报告源级负对数评分、Brier评分和90\%最高后验密度区域覆盖，而非仅报告方差。对离散真源单元$g^*$，
\begin{equation}
 \begin{aligned}
  \mathrm{NLL}_t&=-\log P_t(g^*),\\
  \mathrm{BS}_t&=\sum_{g\in\mathcal G}
      \big[P_t(g)-\I\{g=g^*\}\big]^2.
 \end{aligned}
 \label{eq:evalscores}
\end{equation}
若$P_t(g^*)=0$，NLL应保留无限或明确报告预先登记的数值读出规则；不得靠事后平滑隐藏模型支持失败。错误高置信按预先固定的概率阈值与定位容差定义，报告发生率、持续时间、恢复代价及未恢复比例。官方\texttt{success}保留原定义单独记录，不作为真实位置成功的替代。

同一世界上的方法差采用配对统计，以独立世界而非窗口为重采样单位；跨源/风结论需按其分层报告并给出置信区间。试验规模由先验功效分析或预设置信区间宽度确定，不能依据开发集胜负追加种子。样本效率以相同测试集上的训练世界数量学习曲线评价；成本报告原始样本数、两机总飞行时间、墙钟时间、内存与前向计算次数，区分模拟器成本和算法成本。

研究目标为在非饱和、独立留出测试中，相对同权限强基线实现至少10\%的稳定定位改善，且不降低官方成功率，并改善或保持后验可靠性。该10\%是待验证目标，不是当前结果；正式注册时应固定主指标、相对/绝对定义及基线误差为零时的处理，禁止在Top-1、MAP和Top-5\%之间事后选择最有利项。只有H1获得独立正向证据后，才启动正式M2/M3组合性能验证；失败时保留原结果及结构诊断，不通过更换真源、阈值、航线或测试集挽救成绩。

\subsection{湖岸迁移与当前适用范围}
枫树岭S2是后续应用验证对象。目前完整DAY/NIGHT、多源位置、独立传播实现及局地风输入资格尚未完成，不能把旧工程交互世界当作已验收的湖陆热力气体数据。后续统一采用原生GADEN-RT并记录实际启用的传播、浓度查询和加速模块；“原生”不等于整条仿真已由GPU执行。1\,m与0.5\,m网格的比较应固定物理域、释放、风和查询语义，报告数值敏感性、成本及源判别变化，再决定生产配置。该网格对照是数据资格和资源选择，不单独计为方法创新。

第一篇聚焦双机同步嗅觉观测与源概率定位；更大规模任务分配与多机协同留待后续研究，闭环部署及实飞形成下一阶段验证。气团中心、移动速度、范围和形状可作为独立辅助任务衔接，但其输出需相应观测和标注支撑。本文没有报告完整三维羽流重建、释放率估计或短时态势预测性能。
